\label{sec:method}
\begin{figure*}[t!]
\centering
\includegraphics[width=\textwidth]{figures/overview_v2.pdf}
\caption{\textbf{Overview of \method.} Given RGB-D from sensor(s), we first construct a point cloud of the scene. The point cloud is then used to produce virtual images around the robot workspace. The virtual images are fed to a multi-view transformer model to predict view-specific features, which are then combined to predict action in 3D.}
\label{fig:overview}
\vspace{-3mm}
\end{figure*}
Our goal is to learn a single model that can complete a wide range of manipulation tasks. The input consists of (1) a language description of the task, (2) the current visual state (from RGB-D camera(s)), and (3) the current gripper state (open or closed). The model should predict an action, specified by a target end-effector pose and gripper state at the next \textit{key-frame}. The key-frames represent important or bottleneck steps of the gripper during the task execution~\cite{johns2021coarse}, such as a pre-pick, grasp, or place pose. Given a target end effector pose, we assume a low-level motion planner and controller that can move the end effector to the target pose. To train the model, we assume a dataset $\mathcal{D}=\{D_1, D_2, \cdots, D_n\}$ of $n$ expert demonstrations covering various tasks is given. Each demonstration $D_i = (\{o^i_{1...m_i}\}, \{a^i_{1...m_i}\}, l_i)$ is a successful roll-out of length $m_i$, where $l_i$ is the language description of the task, $\{o^i_1, o^i_2, ... , o^i_{m_i}\}$ is a sequence of the observations from RGB-D camera(s) with gripper state, and $\{a^i_1, a^i_2, ... , a^i_{m_i}\}$ is the sequence of corresponding robot actions. This demonstration dataset can be used to train models with behavior cloning.

Our proposed method (\method) is a transformer model~\cite{vaswani2017attention} that processes images re-rendered around the robot workspace, produces an output for each view, and then back-projects into 3D to predict gripper pose actions, as shown in Fig.~\ref{fig:overview}.

\smallsec{Rendering} The first step is the re-rendering of camera input. Given the RGB-D image(s) captured by one or multiple sensor cameras, we first reconstruct a point cloud of the scene. The point cloud is then re-rendered from a set of virtual viewpoints anchored in the space centered at the robot's base (see Fig.~\ref{fig:overview} and Fig.~\ref{fig:view}). Specifically, for each view, we render three image maps with a total of 7 channels: (1) RGB (3 channels), (2) depth (1 channel), and (3) $(x, y, z)$ coordinates of the points in the world frame (3 channels). The $(x, y, z)$ coordinates help establish the correspondence of pixels across views, i.e., if pixels from different views share the same $(x, y, z)$, they correspond to the same point in 3D. We use PyTorch3D~\cite{ravi2020pytorch3d} for rendering. We empirically verify various design choices in our rendering pipeline (see Tab.~\ref{table:rlbench_ablation} (left)).

The re-rendering process decouples the input images to the ones fed to the transformer. This offers several benefits such as: the ability to re-render at arbitrary and useful locations (e.g., directly above the table) while not being constrained by real-world camera placements; multi-view reasoning even with a single sensor camera; allowing the use of orthographic images instead of generally provided perspective ones; facilitating 3D point-cloud augmentations and enabling additional channels like point correspondence which are not natively presented in the sensor images. We empirically find that these contribute to achieving high performance with view-based networks (see Sec.~\ref{sec:exp_sim}). 
 
\smallsec{Joint Transformer}
The re-rendered images, the language description of the task, and the gripper state (open or close) are processed by a joint transformer model (see Fig.~\ref{fig:transformer} in the appendix). For language, we use pretrained CLIP~\cite{radford2021learning} embeddings (ResNet-50 variant), which provide one token for each word. For the virtual images, we break each of them into $20 \times 20$ patches and pass through a multi-layer perceptron (MLP) to produce image tokens, similar to ViT~\cite{dosovitskiy2020image}. For the gripper state, similar to PerAct~\cite{peract2022arxiv}, we pass it through an MLP and concatenate it to the image tokens. We also add positional embeddings to all the image and language tokens to preserve the positional information.

Overall, \method has eight self-attention layers. In the first four layers, an image token is only allowed to attend to other tokens from the same image. This biases the network to process individual images first before sharing information across images. We concatenate all the image tokens along with the language tokens afterward. In the last four layers, we allow the attention layers to propagate and accumulate information across different images and text. Finally, the image tokens are rearranged back to the original spatial configuration, resulting in the feature channels of each image. 

\smallsec{Action Prediction} The model outputs an 8-dimensional action, including the 6-DoF target end effector pose (3-DoF for translation and 3-DoF for rotation), 1-DoF gripper state (open or close), and a binary indicator for whether to allow collision for the low-level motion planner (see \cite{peract2022arxiv} for details). For translation, we first predict a heatmap for each view from the per-image features from the joint transformer (as shown in Fig.~\ref{fig:transformer} in the appendix). The heatmaps across different views are then back-projected to predict scores for a discretized set of 3D points that densely cover the robot workspace. Finally, the end effector translation is determined by the 3D point with the highest score. Note that this multi-view heatmap representation for translation prediction extends prior approaches in the 2D top-down view setting~\cite{transporter2021}. Hence, \method inherits the benefit of superior sample efficiency by representing the visual input and action in the same spatial structure~\cite{transporter2021}. 


For end effector rotation, we follow PerAct to use the Euler angles representation, where each angle is discretized into bins of $5^{\circ}$ resolution. The gripper state and the motion planner collision indicator are represented as binary variables. To predict the rotations, gripper state, and collision indicator, we use global features ($\mathcal{G}$). The global features are a concatenation of (1) the sum of image features along the spatial dimensions, weighted by the predicted translation heatmap; and (2) the max-pooled image features along the spatial dimension. Specifically, let $f_i$ be the image feature and $h_i$ be the predicted translation heatmap for the $i$th image. Then the global feature $\mathcal{G}$ is given by $\mathcal{G} = \left[ \phi(f_1 \odot h_1); \cdots ;\phi(f_K \odot h_K); \psi(f_1); \cdots; \psi(f_K) \right]$, where $K$ is the number of images, $\odot$ denotes element-wise multiplication, and $\phi$ and $\psi$ denote the sum and max-pooling over the height and width dimensions. The weighted sum operation provides higher weights to image locations near the predicted end effector position.


\smallsec{Loss Function}
We train \method using a mixture of losses. For heatmaps, we use the cross-entropy loss for each image. The ground truth is obtained by a truncated Gaussian distribution around the 2D projection of the ground-truth 3D location. For rotation, we use the cross-entropy loss for each of the Euler angles. We use binary classification loss for the gripper state and collision indicator.